\chapter{Diseño}
\label{chap:diseno}

\lettrine{E}{l} diseño de PadelCenter gira en torno a tres decisiones
arquitectónicas fundamentales: la arquitectura hexagonal como estructura del
backend, el enfoque \textit{contract-first} con OpenAPI como contrato de la API,
y OAuth2 Resource Server con JWT local como mecanismo de seguridad. En este
capítulo se describen estas decisiones y sus consecuencias en la estructura del
sistema.

\section{Arquitectura hexagonal}
\label{sec:arquitectura_hexagonal}

La arquitectura hexagonal organiza el backend en tres capas con dependencias
estrictamente dirigidas hacia el interior:

\begin{description}
  \item[\texttt{domain}:] Contiene los modelos de negocio, los enumerados, los
    puertos (interfaces) y las excepciones de dominio. No tiene ninguna dependencia
    externa: ni Spring, ni JPA, ni ninguna otra librería de infraestructura. Es el
    núcleo del sistema.

  \item[\texttt{application}:] Contiene los casos de uso, organizados en cuatro
    subpaquetes según su naturaleza: \texttt{create}, \texttt{read},
    \texttt{update} y \texttt{delete}. También define los tipos
    \texttt{CommandUseCase} y \texttt{QueryUseCase} como interfaces base. Solo
    depende de \texttt{domain}.

  \item[\texttt{infrastructure}:] Contiene los adaptadores de entrada
    (\texttt{inbound}: controladores HTTP, mappers de API) y de salida
    (\texttt{outbound}: entidades JPA, repositorios, mappers de persistencia,
    seguridad). Implementa los puertos definidos en \texttt{domain}.
\end{description}

\begin{figure}[ht]
\centering
\fbox{\parbox{0.85\textwidth}{
\centering
\vspace{1em}
\textbf{infrastructure.inbound} (controllers, API mappers) \\[0.3em]
\textbf{infrastructure.outbound} (JPA entities, persistence mappers, JWT) \\[0.5em]
$\downarrow$ depende de \\[0.5em]
\textbf{application} (casos de uso: create / read / update / delete) \\[0.5em]
$\downarrow$ depende de \\[0.5em]
\textbf{domain} (records, puertos, excepciones — sin dependencias externas)
\vspace{1em}
}}
\caption{Estructura de capas de la arquitectura hexagonal de PadelCenter.}
\label{fig:arquitectura_hexagonal}
\end{figure}

\subsection{Modelos de dominio como records inmutables}

Una decisión de diseño relevante es que todos los modelos del dominio son
\textbf{Java records} (introducidos en Java 16 y estables desde Java 17). Los
records son clases inmutables por definición: todos sus campos se declaran en el
constructor canónico, no tienen setters y proporcionan \texttt{equals},
\texttt{hashCode} y \texttt{toString} automáticamente.

Esta elección refuerza la arquitectura hexagonal de dos formas:
\begin{itemize}
  \item Ningún adaptador puede modificar un objeto de dominio una vez creado,
    evitando efectos secundarios difíciles de rastrear.
  \item Los records no pueden extender otras clases, lo que impide que las
    entidades de dominio hereden accidentalmente de clases de infraestructura
    (como entidades JPA).
\end{itemize}

Las entidades JPA (\texttt{UserEntity}, \texttt{BookingEntity}, etc.) son clases
convencionales en la capa de infraestructura y no se exponen nunca fuera de ella.

\section{Modelo de dominio}
\label{sec:modelo_dominio}

Las entidades principales del dominio son:

\begin{description}
  \item[\texttt{User}:] Usuario del sistema. Campos: \texttt{userId} (UUID),
    \texttt{firstName}, \texttt{lastName}, \texttt{email}, \texttt{passwordHash}
    (BCrypt), \texttt{phoneNumber}, \texttt{centerRoles}, \texttt{audit}.

  \item[\texttt{Center}:] Centro deportivo. Campos: \texttt{centerId} (UUID),
    \texttt{name}, \texttt{address}, \texttt{city}, \texttt{phoneNumber},
    \texttt{email}, \texttt{manager}, \texttt{audit}.

  \item[\texttt{Field}:] Pista de un centro. Campos: \texttt{fieldId} (UUID),
    \texttt{name}, \texttt{type}, \texttt{pricePerHour} (\texttt{BigDecimal}),
    \texttt{isAvailable} (Boolean), \texttt{center}, \texttt{audit}.

  \item[\texttt{Booking}:] Reserva. Campos: \texttt{bookingId} (\texttt{BIGINT}
    — excepción histórica respecto al uso de UUID en el resto de entidades),
    \texttt{user}, \texttt{field}, \texttt{startTime}/\texttt{endTime}
    (\texttt{LocalDateTime}), \texttt{totalPrice}, \texttt{status}
    (\texttt{PENDING}, \texttt{CONFIRMED}, \texttt{CANCELLED}, \texttt{COMPLETED}),
    \texttt{bookedAt}, \texttt{audit}.

  \item[\texttt{CenterRole}:] Rol contextual de un usuario en un centro concreto.
    Permite que un mismo usuario sea \texttt{ADMIN} en un centro y \texttt{USER}
    en otro.

  \item[\texttt{Tournament}:] Torneo. Campos: \texttt{tournamentId} (UUID),
    \texttt{centerId}, \texttt{name}, \texttt{format}
    (\texttt{ROUND\_ROBIN}, \texttt{ELIMINATION} o
    \texttt{GROUPS\_AND\_ELIMINATION}),
    \texttt{status} (ver figura~\ref{fig:ciclo_torneo}), \texttt{maxPairs},
    \texttt{startDate}, \texttt{endDate}, \texttt{audit}.

  \item[\texttt{TournamentPair}:] Pareja inscrita en un torneo. Estado:
    \texttt{PENDING}, que evoluciona a \texttt{CONFIRMED} o \texttt{REJECTED}.
    Incluye \texttt{teamName} opcional.

  \item[\texttt{Match}:] Partido generado para un emparejamiento. Estado:
    \texttt{SCHEDULED}, \texttt{IN\_PROGRESS}, \texttt{COMPLETED},
    \texttt{CANCELLED}. Contiene \texttt{MatchResult} con marcador y ganador.

  \item[\texttt{Auditable}:] Record base de auditoría incluido en todas las
    entidades: \texttt{createdAt}, \texttt{createdBy} (UUID),
    \texttt{modifiedAt}, \texttt{modifiedBy}, \texttt{deletedAt},
    \texttt{deletedBy} (soft delete).
\end{description}

\section{Diseño de la API — Contract-First}
\label{sec:diseno_api}

El enfoque \textit{contract-first} implica que la especificación OpenAPI se
escribe antes que cualquier código de implementación. El fichero
\texttt{docs/openapi/index.yaml} es la fuente de verdad del contrato de la API.

El proceso de desarrollo de un nuevo endpoint sigue estos pasos:
\begin{enumerate}
  \item Se añade el endpoint al fichero \texttt{index.yaml} con sus parámetros,
    esquemas de request y response, y códigos de estado.
  \item Se ejecuta \texttt{./mvnw generate-sources}, que invoca el plugin
    \textit{openapi-generator-maven-plugin~7.6.0} y genera las interfaces de
    controlador y los DTOs (como records Java) en \texttt{target/generated-sources}.
  \item El controlador implementa la interfaz generada. Si la firma del método no
    coincide con la especificación, el compilador lo detecta.
  \item En el frontend, \texttt{pnpm generate-api} (con el backend en
    \texttt{:8080}) regenera los hooks React Query mediante Orval.
\end{enumerate}

\subsection{Convenciones de la especificación}

\begin{itemize}
  \item Todos los endpoints siguen el prefijo \texttt{/api/v1/}.
  \item Los identificadores de recursos son UUIDs (salvo \texttt{bookingId},
    que es BIGINT por razones históricas).
  \item Los errores siguen el esquema \texttt{ProblemDetail} de RFC~9457.
  \item Los \textit{timestamps} se representan en formato ISO~8601 con zona
    horaria UTC.
  \item La paginación sigue el esquema de Spring Data con \texttt{page},
    \texttt{size} y \texttt{sort}, envuelto en \texttt{PageResult<T>}.
\end{itemize}

\subsection{Endpoints de la API}

La tabla~\ref{tab:endpoints} recoge los 40 endpoints expuestos por la API,
organizados por controlador.

\begin{table}[ht]
\centering
\resizebox{\textwidth}{!}{%
\begin{tabular}{llll}
\hline
\textbf{Método} & \textbf{Ruta} & \textbf{Descripción} & \textbf{Rol mínimo} \\
\hline
POST & /api/v1/auth/login & Iniciar sesión & Anónimo \\
POST & /api/v1/auth/refresh & Renovar token & Anónimo \\
POST & /api/v1/auth/logout & Cerrar sesión & Anónimo \\
GET  & /api/v1/me & Usuario actual & USER \\
GET  & /api/v1/me/bookings & Mis reservas & USER \\
GET  & /api/v1/me/tournaments & Mis torneos & USER \\
GET  & /api/v1/me/matches & Mis partidos & USER \\
\hline
POST & /api/v1/users & Registrarse & Anónimo \\
GET  & /api/v1/users & Listar usuarios & ADMIN \\
GET  & /api/v1/users/\{id\} & Ver usuario & ADMIN \\
PUT  & /api/v1/users/\{id\} & Actualizar usuario & Propietario/ADMIN \\
PUT  & /api/v1/users/\{id\}/roles & Cambiar roles & ADMIN \\
PUT  & /api/v1/users/\{id\}/password & Cambiar contraseña & Propietario \\
DELETE & /api/v1/users/\{id\} & Eliminar usuario & ADMIN \\
GET  & /api/v1/users/\{id\}/bookings & Reservas de usuario & ADMIN \\
\hline
POST & /api/v1/centers & Crear centro & ADMIN \\
GET  & /api/v1/centers & Listar centros & USER \\
GET  & /api/v1/centers/\{id\} & Ver centro & USER \\
DELETE & /api/v1/centers/\{id\} & Eliminar centro & ADMIN \\
POST & /api/v1/centers/\{id\}/fields & Crear pista & ADMIN \\
GET  & /api/v1/centers/\{id\}/fields & Listar pistas del centro & USER \\
GET  & /api/v1/centers/\{id\}/bookings & Reservas del centro & MANAGER \\
GET  & /api/v1/centers/\{id\}/tournaments & Torneos del centro & USER \\
\hline
GET  & /api/v1/fields & Listar pistas & USER \\
GET  & /api/v1/fields/\{id\} & Ver pista & USER \\
GET  & /api/v1/fields/\{id\}/availability & Disponibilidad por fecha & USER \\
PUT  & /api/v1/fields/\{id\}/availability & Activar/desactivar pista & MANAGER \\
DELETE & /api/v1/fields/\{id\} & Eliminar pista & ADMIN \\
\hline
POST & /api/v1/bookings & Crear reserva & USER \\
GET  & /api/v1/bookings/\{id\} & Ver reserva & Propietario/MANAGER \\
PUT  & /api/v1/bookings/\{id\} & Actualizar reserva & Propietario \\
POST & /api/v1/bookings/\{id\}/cancel & Cancelar reserva & Propietario/MANAGER \\
\hline
POST & /api/v1/tournaments & Crear torneo & ADMIN \\
GET  & /api/v1/tournaments/\{id\} & Ver torneo & USER \\
DELETE & /api/v1/tournaments/\{id\} & Eliminar torneo & ADMIN \\
POST & /api/v1/tournaments/\{id\}/pairs & Inscribirse como pareja & USER \\
GET  & /api/v1/tournaments/\{id\}/pairs & Ver parejas & USER \\
POST & /api/v1/tournaments/\{id\}/pairs/\{pId\}/confirm & Confirmar pareja & MANAGER \\
POST & /api/v1/tournaments/\{id\}/registration/open & Abrir inscripción & ADMIN \\
POST & /api/v1/tournaments/\{id\}/registration/close & Cerrar e generar cuadro & ADMIN \\
GET  & /api/v1/tournaments/\{id\}/matches & Ver partidos & USER \\
GET  & /api/v1/tournaments/\{id\}/standings & Ver clasificación & USER \\
\hline
POST & /api/v1/matches/\{id\}/result & Reportar resultado & MANAGER \\
\hline
\end{tabular}}
\caption{Endpoints de la API REST de PadelCenter.}
\label{tab:endpoints}
\end{table}

\section{Esquema de base de datos}
\label{sec:esquema_bd}

La base de datos es PostgreSQL~16~\cite{postgresql}. Las migraciones se gestionan con Flyway y se
nombran siguiendo el convenio \texttt{V\{version\}\_\_\{descripcion\}.sql}.

Las tablas principales son:

\begin{itemize}
  \item \texttt{users}: datos de usuario y credenciales (BCrypt).
  \item \texttt{center\_roles}: relación usuario-centro-rol con restricción
    \texttt{UNIQUE(user\_id, center\_id)}.
  \item \texttt{centers}: centros deportivos con soft delete.
  \item \texttt{fields}: pistas de cada centro con precio y flag
    \texttt{is\_available}.
  \item \texttt{bookings}: reservas con clave primaria \texttt{BIGSERIAL} y
    control de solapamiento mediante índice parcial (migración V4).
  \item \texttt{booking\_participants}: participantes adicionales de una reserva.
  \item \texttt{tournaments}: torneos con \texttt{CHECK} sobre \texttt{status}
    (6 valores) y \texttt{format} (3 valores).
  \item \texttt{tournament\_pairs}: parejas inscritas con \texttt{team\_name}
    (añadido en migración V8).
  \item \texttt{tournament\_matches}: partidos con resultado embebido.
\end{itemize}

Todas las tablas incluyen columnas de auditoría: \texttt{created\_at},
\texttt{created\_by} (UUID), \texttt{modified\_at}, \texttt{modified\_by},
\texttt{deleted\_at}, \texttt{deleted\_by}.

\section{Diseño de seguridad}
\label{sec:diseno_seguridad}

La autenticación se basa en JWT firmados con HMAC-SHA256 (\texttt{HS256}).
En el lado del servidor, la verificación se implementa mediante el módulo
{\ttfamily spring-\-boot-\-starter-\-oauth2-\-resource-\-server} de Spring Security, que
configura un \texttt{NimbusJwtDecoder} con la clave secreta local. Esto permite
aprovechar toda la infraestructura de integración JWT de Spring
(decodificación, validación de claims, inyección en el contexto de seguridad)
sin depender de ningún proveedor de identidad externo.

El flujo de autenticación es el siguiente:
\begin{enumerate}
  \item El cliente envía \texttt{POST /api/v1/auth/login} con correo y contraseña.
  \item El servidor verifica las credenciales y genera un \textit{access token}
    (15 minutos de validez) y un \textit{refresh token} (7 días).
  \item El \textit{access token} se devuelve en el cuerpo de la respuesta; el
    \textit{refresh token} se establece como cookie \textit{httpOnly} y
    \texttt{SameSite=Strict}.
  \item Las peticiones protegidas incluyen el \textit{access token} en el
    header \texttt{Authorization: Bearer \{token\}}.
  \item Cuando el \textit{access token} expira, el cliente llama a
    \texttt{POST /api/v1/auth/refresh}; el servidor lee la cookie y emite un
    nuevo par de tokens.
\end{enumerate}

\subsection{Evaluadores de autorización personalizados}

Más allá de la autenticación, PadelCenter implementa autorización granular
mediante evaluadores custom registrados como beans de Spring:

\begin{description}
  \item[\texttt{@bookingOwnerEvaluator}:] Verifica que el usuario autenticado
    es el propietario de la reserva, o que tiene rol \texttt{MANAGER} o
    \texttt{ADMIN} en el centro de la pista reservada.

  \item[\texttt{@centerRoleEvaluator}:] Verifica que el usuario tiene el rol
    mínimo requerido en el centro sobre el que opera. Permite operaciones
    diferenciadas según si el actor es \texttt{USER}, \texttt{MANAGER} o
    \texttt{ADMIN} en ese centro concreto.

  \item[\texttt{@tournamentRoleEvaluator}:] Verifica permisos específicos
    sobre torneos: inscripción (cualquier \texttt{USER}), confirmación de parejas
    y reporte de resultados (\texttt{MANAGER} del centro), gestión del ciclo
    de vida (\texttt{ADMIN}).
\end{description}

Estos evaluadores se invocan mediante \texttt{@PreAuthorize} directamente en
los métodos de los casos de uso, manteniendo la lógica de autorización cerca
de la lógica de negocio sin contaminar los controladores.
